\documentclass[11pt]{article}

\usepackage[margin=1in]{geometry}
\usepackage{amsmath,amssymb,amsthm}
\usepackage{booktabs}
\usepackage{enumitem}
\usepackage{flafter}
\usepackage{graphicx}
\usepackage{placeins}
\usepackage[hidelinks]{hyperref}

\setlist[itemize]{leftmargin=1.4em}
\setlist[enumerate]{leftmargin=1.6em}

\newcommand{\KL}{D}
\newcommand{\E}{\mathbb{E}}
\newcommand{\Nseq}{N}
\newcommand{\model}{\mathcal{M}_{d,L}}
\newcommand{\cstar}{c^\star}
\newtheorem{proposition}{Proposition}
\newtheorem{remark}{Remark}

\title{A Product Model for Sources with Memory}
\author{}
\date{\today}

\begin{document}
\maketitle

\begin{abstract}
The layered product-simplex predictor of~\cite{layered2026} is extended
from exchangeable (memoryless) sources to sources with memory.  The past
is summarized in a state and the scheme predicts from (state, next
token) pairs; the simplest instance is a first-order Markov model whose
state is the previous token.  Every number reported here is the exact
codelength of a complete scheme, produced by a script in the
repository.
\end{abstract}

\section*{The scoreboard}
\label{sec:scoreboard}

Every number in this paper is the honest codelength of a complete
scheme --- the whole file reproduced exactly, every model choice and
the vocabulary paid for --- in bits per character of the original
file, the same accounting under which the published records are set.
This table is the state of the program; each section either moves a
number here or says plainly that it does not.

\paragraph{One accounting convention throughout.}
Every rate is total coded bits divided by the number of characters (bytes) in
the original corpus, never by the number of derived tokens.  When a run uses
an external tokenizer, its vocabulary description is charged.  When it keeps
only $V-1$ token types and maps the remainder to an escape symbol, the selected
subset is transmitted, the reduced stream is coded by the stated predictor,
and each original token hidden by escape is coded by a sequential
Jeffreys--KT code over the excluded tokenizer identifiers.  Any initial symbols
for which the stated memory or checkpoint predictor is unavailable are coded
by the declared lower-order or fixed-alphabet KT rule.  The sum of these
components is therefore a lossless code for the original file.  Sections below
state only what differs from this convention: the token representation,
retained alphabet, state or context, and predictive rule.

\begin{table}[ht]
\centering
\begin{tabular}{lccc}
\toprule
 & \texttt{text8} & \texttt{enwik8} & \texttt{enwik9} \\
\midrule
bytes, no memory                & $4.1235$ & $5.0802$ & $5.1565$ \\
LLM tokenizer, no memory        & $2.1716$ & $2.9552$ & $2.9775$ \\
first order (Sec.~\ref{sec:firstorder})   & $1.8298$ & $2.1135$ & $1.8371$ \\
second-order Markov (Sec.~\ref{sec:ordertwo}) & $1.8298$ & $2.1135$ & $1.7915$ \\
context tree (Sec.~\ref{sec:ctree})       & $\mathbf{1.8182}$ & $\mathbf{2.0606}$ & $\mathbf{1.6598}$ \\
calibrated order two (Sec.~\ref{sec:calibration}) & \multicolumn{3}{c}{controlled rerun in progress} \\
\midrule
best published (Table~\ref{tab:files})    & $1.398$  & $1.170$  & $0.858$ \\
\bottomrule
\end{tabular}
\caption*{Best honest bits per character achieved by each
  construction of this paper, against the published records
  (decompressor size included, no external data --- the same rules).
  Bold: the current bests.}
\end{table}

\section{The model}
\label{sec:setup}

Throughout, the data are a fixed file: a sequence
$x^n = (x_1, \dots, x_n)$ of $n$ tokens from an alphabet of size $d$,
with counts $m_i = \#\{t : x_t = i\}$.  Nothing is assumed about where
the file came from.

\paragraph{A code is a probability assignment.}  If $Q$ gives every
length-$n$ sequence a nonnegative probability and these sum to one,
then arithmetic coding turns $Q$ into a lossless code that spends
$-\log_2 Q(x^n)$ bits on the file $x^n$, up to $O(1/n)$; conversely,
the lengths of any lossless code define such a $Q$.  Choosing a model
means choosing $Q$: a rule, fixed in advance and known to the
decoder, that assigns a probability to every candidate file.

The rule used throughout is the layered product-simplex construction
of~\cite{layered2026}.  It is built in three steps.

\emph{A weighting over the simplex.}  Fix a depth $L \ge 1$.  Draw $L$
points $\theta^{(1)}, \dots, \theta^{(L)}$ of the $d$-simplex
independently and uniformly, and form the coordinatewise product,
renormalized:
\begin{equation}
\label{eq:product-prior}
  \theta_i \;=\;
  \frac{\prod_{\ell=1}^{L} \theta^{(\ell)}_i}
       {\sum_{j=1}^{d} \prod_{\ell=1}^{L} \theta^{(\ell)}_j},
  \qquad i = 1, \dots, d.
\end{equation}
The distribution of the resulting vector $\theta$ is a measure
$\pi_L$ on the simplex.  At $L=1$ it is uniform; as $L$ grows it
concentrates on increasingly sparse, low-entropy vectors.

\emph{The assignment at depth $L$.}  For every finite string
$x^t = (x_1, \dots, x_t)$, of any length, define
\begin{equation}
\label{eq:mixture}
  Q^{(L)}(x^t)
  \;=\;
  \int \prod_{s=1}^{t} \theta_{x_s} \; \pi_L(d\theta),
\end{equation}
the integral being over the vector $\theta$ with respect to $\pi_L$.
Restricted to the $d^t$ strings of any one length $t$, these values
sum to one: summing the product over strings inside the integral
gives $(\sum_i \theta_i)^t = 1$.  So \eqref{eq:mixture} is a
probability assignment at every length.  Since the product depends on
the string only through its counts, so does $Q^{(L)}$: the assignment
is \emph{exchangeable}.

\emph{The average over depths.}  The depth-averaged assignment is
\begin{equation}
\label{eq:depth-average}
  Q^{\mathrm{avg}}(x^t)
  \;=\;
  \frac{1}{L_{\max}}\sum_{L=1}^{L_{\max}} Q^{(L)}(x^t),
  \qquad
  L_{\max}=\operatorname{round}(2\cstar\ln d),
  \quad
  \cstar=\frac{1}{1-\gamma}\approx 2.365,
\end{equation}
again defined on every finite string.

\paragraph{Sequential decoding.}  A decoder reconstructs the file one
symbol at a time: at step $t$ it needs the probability of each
candidate next symbol, computed from the symbols it has already
recovered.

\emph{Consistency across lengths.}  For every prefix $x^t$,
\begin{equation}
\label{eq:consistency}
  \sum_{y} Q\bigl(x^{t}y\bigr) \;=\; Q\bigl(x^{t}\bigr)
\end{equation}
holds for $Q = Q^{(L)}$, and hence for $Q = Q^{\mathrm{avg}}$.  For
$Q^{(L)}$, sum over $y$ inside the integral \eqref{eq:mixture} and
use $\sum_y \theta_y = 1$.  For $Q^{\mathrm{avg}}$, sum the identity
for $Q^{(L)}$ over $L$ with weights $1/L_{\max}$.  Consequently
\begin{equation*}
  Q(y \mid x^{t}) \;=\; \frac{Q(x^{t}y)}{Q(x^{t})}
\end{equation*}
is a probability distribution over $y$ at every step, and both
numerator and denominator depend only on symbols already decoded.

\emph{The conditional.}  At a fixed depth, the ratio
$Q^{(L)}(x^t y)/Q^{(L)}(x^t)$ can be written as an integral against a
new measure on the simplex:
\begin{equation}
\label{eq:posterior-mean}
  Q^{(L)}(y \mid x^{t})
  \;=\; \int \theta_y \; \rho_t(d\theta),
  \qquad
  \rho_t(d\theta)
  \;=\;
  \frac{\prod_{s=1}^{t} \theta_{x_s}\; \pi_L(d\theta)}
       {\displaystyle\int
        \prod_{s=1}^{t} \tilde\theta_{x_s}\; \pi_L(d\tilde\theta)}.
\end{equation}
At $L = 1$ the integrals
are explicit and \eqref{eq:posterior-mean} equals
$(m_y + 1)/(t + d)$, Laplace's rule, with $m_y$ the number of
occurrences of $y$ in $x^t$.

For the depth average, define
\begin{equation}
\label{eq:depth-posterior}
  w_L(x^t) \;=\;
  \frac{Q^{(L)}(x^t)}{\sum_{L'=1}^{L_{\max}} Q^{(L')}(x^t)}.
\end{equation}
Expanding the ratio,
\begin{equation*}
  Q^{\mathrm{avg}}(y \mid x^{t})
  \;=\; \sum_{L=1}^{L_{\max}} w_L(x^{t})\, Q^{(L)}(y \mid x^{t}).
\end{equation*}
The prediction at step $t$ is the average of the per-depth
predictions \eqref{eq:posterior-mean} with the weights
\eqref{eq:depth-posterior}, not with the uniform weights
$1/L_{\max}$; the weights \eqref{eq:depth-posterior} change with
every decoded symbol.

\emph{The codelength needs one evaluation.}  Encoder and decoder run
the conditionals step by step.  The codelength they produce does not
have to be accumulated that way: consecutive ratios telescope, so
\begin{equation*}
  -\sum_{t=1}^{n} \log_2 Q^{\mathrm{avg}}(x_t \mid x^{t-1})
  \;=\; -\log_2 Q^{\mathrm{avg}}(x^n),
\end{equation*}
one evaluation at the final counts.  By exchangeability this depends
on $x^n$ only through its count profile $\lambda$ (the multiset of
nonzero counts):
$-\log_2 Q^{\mathrm{avg}}(x^n) =
-\log_2\bigl[\tfrac1{L_{\max}}\sum_L q_\lambda(L)\bigr]$, where
$q_\lambda(L)$ is evaluated by the layer recursion and saddlepoint
method of Appendix~B of~\cite{layered2026}, ported to this repository
(\texttt{src/product\_model\_with\_memory/layered.py}, with the
batched table builder of \texttt{fast\_tables.py} validated against
the original).  Appendix~\ref{app:factorisation} gives the same
argument per state for the models with memory.

\paragraph{The two quantities we report.}

\emph{Empirical entropy.}  The empirical unigram entropy of the file
is the entropy of its relative frequencies,
\begin{equation}
\label{eq:empirical-entropy}
  \hat H_n \;=\; -\sum_{i\,:\,m_i>0} \frac{m_i}{n}\,
    \log_2 \frac{m_i}{n}
  \qquad \text{[bits/token]}.
\end{equation}
This number is computed from the file itself, so it is a benchmark,
not the length of any code.  It is a floor in the following sense.  A
coder that uses one fixed probability assignment $p$ per symbol for
the whole file spends $-\tfrac1n\sum_t \log_2 p(x_t)$ bits per token,
and this is smallest exactly when $p_i = m_i/n$, the frequencies of
the file itself, where it equals $\hat H_n$.  So no fixed per-symbol
assignment can beat $\hat H_n$.  But no decoder can \emph{achieve} it
either: the decoder cannot know the frequencies of data it has not
yet received.  A real scheme closes this gap in one of two ways.
Either send the frequencies first and then code the data under them,
paying for their description; or use at each step only probabilities
computable from the symbols already decoded, as every scheme in this
paper does, paying instead through the early predictions, which are
made from few observations.  The excess over $\hat H_n$ is the price
of that gap: the cost of learning the distribution while coding under
it.

\emph{Loss and redundancy.}  The loss of an assignment $Q$ on the
file is its per-token codelength
\begin{equation}
\label{eq:codelength}
  \ell_n(Q) \;=\; -\tfrac{1}{n}\log_2 Q(x^n)
  \;=\; -\tfrac{1}{n}\sum_{t=1}^{n}
        \log_2 Q\bigl(x_t \,\big|\, x^{t-1}\bigr),
\end{equation}
the two forms equal by the chain rule, and realizable as an actual
arithmetic code by the preceding paragraphs.  The \emph{redundancy}
we tabulate is $\ell_n(Q) - \hat H_n$: the part of the codelength
paid for \emph{learning} the distribution rather than describing the
data under a known one.

\emph{Depth diagnostic.}  We also report the mode
$L^\star = \arg\max_L w_L(x^n)$ of the depth weights
\eqref{eq:depth-posterior} at the final counts.  The mode plays no
role in any codelength, which is always that of the full average
\eqref{eq:depth-average}; it is a diagnostic that names the single
depth best explaining the data, and its position relative to
$\ln d$ locates the source on the complexity spectrum
of~\cite{layered2026}.

\section{Sources and representations}
\label{sec:sources}

\subsection*{Files}

Every result in this paper is measured on the three files of the
standard text-compression benchmarks (Table~\ref{tab:files}).
\texttt{text8} is the first $10^8$ bytes of an English Wikipedia dump
with everything removed except lowercase \texttt{a}--\texttt{z} and the
space: no markup, no capitals, no punctuation, no digits.
\texttt{enwik8} and \texttt{enwik9} are the first $10^8$ and $10^9$
bytes of the same dump left exactly as it is --- XML tags, mixed case,
punctuation, digits, and text in many languages.  The three differ by roughly one bit per character at every
stage below.

\begin{table}[ht]
\centering
\begin{tabular}{lrrll}
\toprule
file & bytes & byte values & contents & best published \\
\midrule
\texttt{text8}  & $10^8$ & $27$  & cleaned lowercase text & $1.398$ \\
\texttt{enwik8} & $10^8$ & $205$ & raw Wikipedia XML      & $1.170$ \\
\texttt{enwik9} & $10^9$ & $206$ & raw Wikipedia XML      & $0.858$ \\
\bottomrule
\end{tabular}
\caption{The three files.  ``Byte values'' is how many of the $256$
  possible bytes actually occur.  ``Best published'' is the best result
  on the public lists, in bits per character, by \texttt{paq8h},
  \texttt{cmix~v21} and \texttt{nncp~v3.2} respectively.  Those programs
  use memory heavily, so they are the target of this work rather than
  a comparison for its early sections.  \texttt{text8} md5
  \texttt{3bea1919949baf155f99411df5fada7e}.}
\label{tab:files}
\end{table}

\paragraph{Bits per character means bits per byte.}  Every figure in
this paper is compressed bits divided by the file's \emph{byte}
count, $10^{8}$ or $10^{9}$.  On \texttt{text8} bytes and characters
coincide exactly (the file is plain ASCII).  On \texttt{enwik8} and
\texttt{enwik9} they differ slightly --- the files are UTF-8, and
some characters occupy several bytes --- so ``bits per character''
is, strictly, bits per byte there.  We keep the name because the
compression literature has always used it, and the published records
of Table~\ref{tab:files} are computed with the same divisor, so
every comparison in this paper is like against like.

\subsection*{Symbols}

Before anything can be compressed we have to say what a symbol is.  This
paper uses two answers, and everything else follows from them.

\paragraph{Bytes.}  Every byte is a symbol, so the alphabet has $256$
letters and never changes.  Nothing has to be explained to the decoder
and bits per symbol are bits per character directly.  This is the honest
floor: no tokenizer, no assumptions, and any scheme that cannot beat it
is doing nothing.

\paragraph{The LLM tokenizer.}  Language models are trained on a fixed
subword vocabulary, built in advance from a large body of text, and
every large model of recent years sees its training data as a stream of
such tokens.  We use \texttt{cl100k\_base} --- the vocabulary behind
GPT-3.5 and GPT-4 --- through the \texttt{tiktoken} library:
$100{,}277$ subwords, identical for every file.  This is the
representation that matters here, because the long-term aim of this
project is a language model built on the same estimator, and a result
measured on any other token stream would not transfer to it.

It has one cost that bytes do not.  The vocabulary comes from outside
the file, so a decompressor must carry it and have its size counted:
$776$ kB zipped, which is $0.0621$ bits per character on a $10^8$-byte
file and $0.0062$ on $10^9$.  We report every figure both ways --- with
that charge, which is the honest number, and without it, which is the
convention of the language-model literature and is not a valid entry.

\begin{table}[ht]
\centering
\begin{tabular}{lll}
\toprule
 & bytes & LLM tokenizer \\
\midrule
one symbol is         & one byte          & a subword from a fixed list \\
alphabet size         & $256$             & $100{,}277$ \\
vocabulary comes from & nothing to learn  & text outside the file \\
must be shipped       & nothing           & $776$ kB zipped \\
valid benchmark entry & yes               & only if that is counted \\
\bottomrule
\end{tabular}
\caption{The two representations, before any measurement.  The last two
  rows decide admissibility: a decompressor may use unlimited
  computation, but everything it carries must be counted, and a
  pretrained vocabulary is such a thing.}
\label{tab:representations}
\end{table}

\section{No memory}
\label{sec:text8}

This section fixes the floor: what the construction achieves with no
memory at all (Table~\ref{tab:memoryless}).  Each file is a single count
profile, so these are one evaluation apiece, and every figure is the
exact codelength of a complete scheme converted to bits per character of
the original file.

\begin{table}[ht]
\centering
\begin{tabular}{lccc}
\toprule
 & \texttt{text8} & \texttt{enwik8} & \texttt{enwik9} \\
\midrule
bytes                                  & $4.1235$ & $5.0802$ & $5.1565$ \\
LLM tokenizer, vocabulary counted      & $2.1716$ & $2.9552$ & $2.9775$ \\
LLM tokenizer, not counted$^{\dagger}$ & $2.1095$ & $2.8931$ & $2.9713$ \\
\bottomrule
\end{tabular}
\caption{Memoryless coding, bits per character of the original file.
  $^{\dagger}$~The last row omits the $776$ kB vocabulary and is
  therefore not a valid entry; it is given so that these results can be
  compared with language-model work on equal terms.}
\label{tab:memoryless}
\end{table}

\paragraph{The estimator costs nothing where there is nothing to learn.}
Over bytes the layered mixture costs $4.123534$ against a plain order-0
entropy of $4.123527$ on \texttt{text8}, and the gap is in the fifth
decimal on all three files ($7\cdot10^{-6}$, $2\cdot10^{-5}$,
$3\cdot10^{-6}$ bits per character).  The same holds on the large alphabet: over the $100{,}277$ subwords of the LLM
tokenizer the redundancy over the plug-in entropy is $0.0150$ bits per
token on \texttt{enwik8} and $0.0020$ on \texttt{enwik9}.  A hundred thousand symbols with tens of millions of
observations is still the dense regime, so every gain reported below is
attributable to memory rather than to the estimator.

\paragraph{Choosing symbols well is worth about two bits per character.}
Coding \texttt{enwik8} byte by byte costs $5.0802$ bits per character,
while coding it as subwords costs $2.9552$ with the vocabulary counted.
The difference is $2.13$ bits per character, obtained with no
context.  It exceeds every gain from memory reported below.

\paragraph{The shipped vocabulary amortizes.}  Its charge falls from
$0.0621$ bits per character at $10^8$ to $0.0062$ at $10^9$; the
counted and uncounted rows of Table~\ref{tab:memoryless} differ only
in the fourth decimal on \texttt{enwik9}.

\section{Memory of order one}
\label{sec:firstorder}

Each symbol is conditioned on the one immediately before it and on
nothing else (Table~\ref{tab:firstorder}).  The state is the previous
symbol, every distinct symbol has its own state, nothing is pooled and
nothing is capped, and each state's successors are coded by the
per-state layered mixture used throughout.  The state map is a function
of the past that the decoder computes for itself, and the product over
states is exactly the codelength of a symbol-by-symbol predictor: the
one that weights the depths, at every step, by their posterior given
what that state has seen so far.  Appendix~\ref{app:factorisation}
proves this, distinguishes it from the receiver that weights the depths
uniformly --- a different code, and a worse one --- and reports the
check that the sequential walk and the closed form agree to $0.000$
bits.  With a single state the same machinery is the
memoryless model, and it reproduces the
Table~\ref{tab:memoryless} figures to four decimals.

\begin{table}[ht]
\centering
\begin{tabular}{lrrrrr}
\toprule
file & memoryless & first order & saved & states & best possible \\
\midrule
\texttt{text8}  & $2.1716$ & $1.8298$ & $0.3418$ & $35{,}767$ & $0.701$ \\
\texttt{enwik8} & $2.9552$ & $2.1135$ & $0.8417$ & $71{,}161$ & $1.242$ \\
\texttt{enwik9} & $2.9775$ & $1.8371$ & $1.1404$ & $81{,}276$ & $1.281$ \\
\bottomrule
\end{tabular}
\caption{First-order memory over the LLM tokenizer, bits per character
  of the original file, vocabulary counted throughout.  ``States'' is
  how many distinct previous symbols occur, so it is the number of
  separate distributions the model must learn.  ``Best possible'' is the
  empirical conditional entropy --- what a first-order model would save
  if handed the true distributions --- so the gap between it and
  ``saved'' is the price of having to learn them from the file.}
\label{tab:firstorder}
\end{table}

\paragraph{Both accounting conventions.}  The figures above count the
vocabulary.  The uncounted convention of Table~\ref{tab:memoryless} is
obtained by subtracting the charge, $0.0621$ bits per character at
$10^8$ and $0.0062$ at $10^9$: $1.7677$ on \texttt{text8}, $2.0514$ on
\texttt{enwik8}, $1.8309$ on \texttt{enwik9}.  The charge does not
depend on the model, so ``saved'' and ``best possible'' are the same
under either convention and only the level changes.

\paragraph{The cost of learning is large.}  In
Section~\ref{sec:text8} one distribution is estimated from $25.8$
million observations and the estimator is free.  First order estimates
$71{,}161$ distributions from the same data, about $362$ observations
each, and captures $68\%$ of what is available on \texttt{enwik8}.  The
remainder is the cost of estimating that many distributions from that
little data.

\paragraph{The cost falls with data per state.}  \texttt{enwik8} and
\texttt{enwik9} isolate this, since the vocabulary is fixed and the
state count barely moves --- $71{,}161$ against $81{,}276$ --- while
the data grows tenfold, from $362$ observations per state to $3{,}368$.
The captured share rises from $68\%$ to $89\%$, and the shortfall falls
from $0.400$ to $0.141$ bits per character.  \texttt{enwik9} is
therefore cheaper per character than \texttt{enwik8} under first-order
memory, $1.8371$ against $2.1135$, although it is the more expensive
file without memory.  The learning cost is a property of the ratio of
data to states, not of the alphabet.

\paragraph{Coarsening the state does not help.}  If the shortfall in
Table~\ref{tab:firstorder} is the price of learning too many
distributions, the model should be able to lower it
by learning fewer.  Fix a grid of values $M$ and let member $M$ give
its own state to $M$ symbols and pool every other symbol into one
shared backoff state, so member $M$ has at most $M+1$ states; $M=0$ is
the memoryless model and $M$ at least the vocabulary size is first
order.  The family is averaged uniformly, which costs $\log_2$ of the
number of members over the whole file --- $3.5$ bits for the eleven
members used here --- and guarantees a codelength within that of the
best member.  Which $M$ symbols keep their own state is part of the code, so the
decoder has to know it before decoding begins.  Ranking symbols by
their frequency in the file being compressed does not satisfy this:
the ranking would have to be transmitted, about $1.5$ million bits at
$d = 100{,}277$, or $0.19$ bits per character on \texttt{enwik8},
which exceeds the entire first-order saving.  Member $M$ therefore
keeps the $M$ smallest vocabulary indices.  The tokenizer assigns
indices in merge order, so a small index is a subword that was
frequent in the tokenizer's training corpus, and the ordering is part
of the vocabulary that Table~\ref{tab:firstorder} already pays for.
At $M=0$ and at full memory the two rules coincide, so
Tables~\ref{tab:memoryless} and~\ref{tab:firstorder} are unaffected.

\begin{table}[ht]
\centering
\begin{tabular}{lrrr}
\toprule
states kept & \texttt{text8} & \texttt{enwik8} & \texttt{enwik9} \\
\midrule
$0$        & $2.1716$ & $2.9552$ & $2.9775$ \\
$1{,}024$  & $1.9805$ & $2.5686$ & $2.4670$ \\
$8{,}192$  & $1.8698$ & $2.3311$ & $2.1686$ \\
$32{,}768$ & $1.8318$ & $2.1666$ & $1.9323$ \\
all        & $1.8298$ & $2.1135$ & $1.8371$ \\
\bottomrule
\end{tabular}
\caption{Coarsening the first-order state, bits per character,
  vocabulary counted.  ``All'' is $35{,}767$, $71{,}161$ and $81{,}276$
  states respectively.  The full grids run to eleven members; the rows
  shown are those the three grids have in common.}
\label{tab:coarsen}
\end{table}

The codelength falls monotonically in $M$ on all three files
(Table~\ref{tab:coarsen}), the posterior places all its weight on the
full model, and the mixture equals the best member to four decimals.
Coarsening never pays.  Each doubling of $M$ still buys between $0.048$
and $0.11$ bits per character on \texttt{enwik8} up to $65{,}536$
states, where the remaining states hold fewer than $400$ observations
each.
Giving a rare symbol its own state costs the layered estimator very
little, while pooling it with seventy thousand others discards nearly
everything it carries, and the family offers only that choice.

The last step is the informative one.  Going from $65{,}536$ states to
all $71{,}161$ gains $0.0049$ bits per character on \texttt{enwik8},
and the corresponding last step on \texttt{text8} ($32{,}768$ to all
$35{,}767$) gains $0.0020$, against tens of thousandths for every
earlier doubling.  First-order memory sits at the point where additional state
stops earning its cost, which is a statement about the data rather than
about the estimator.

\paragraph{Where the excess is.}  The shortfall is therefore not
caused by having too many states, and its composition can be measured
directly.  For each state, compare what the code pays with what a
decoder that already knew that state's successor distribution would
pay; the difference, summed, is the cost of learning.  The empirical
conditional entropy is optimistic by an amount that grows with the
number of states, so it is reported with the Miller--Madow correction
$(\text{support}_s - 1)/2n_s$ nats per state.

\begin{table}[ht]
\centering
\begin{tabular}{lrrrr}
\toprule
file & model & $H(X \mid \text{prev})$ & excess & obs.\ per state \\
\midrule
\texttt{text8}  & $1.8298$ & $1.4943$ & $0.3355$ & $543$ \\
\texttt{enwik8} & $2.1135$ & $1.7359$ & $0.3776$ & $362$ \\
\texttt{enwik9} & $1.8371$ & $1.7062$ & $0.1309$ & $3{,}367$ \\
\bottomrule
\end{tabular}
\caption{The cost of learning under first-order memory, bits per
  character of the original file, vocabulary counted, on the scale of
  every other table here.  The conditional entropy carries the
  Miller--Madow correction of Appendix~\ref{app:mm}, which is $0.026$,
  $0.027$ and $0.011$ bits per character.  ``Excess'' is the difference
  between the two columns to its left; the vocabulary charge is common
  to both and cancels.}
\label{tab:excess}
\end{table}

\begin{table}[ht]
\centering
\begin{tabular}{lrrr}
\toprule
observations in the state & \texttt{text8} & \texttt{enwik8} & \texttt{enwik9} \\
\midrule
$1$ to $29$          & $3.6\%$  & $7.5\%$  & $1.1\%$ \\
$30$ to $99$         & $10.1\%$ & $14.9\%$ & $2.6\%$ \\
$100$ to $999$       & $41.3\%$ & $41.8\%$ & $25.6\%$ \\
$1{,}000$ to $9{,}999$ & $33.2\%$ & $22.9\%$ & $42.3\%$ \\
$10{,}000$ and above & $11.8\%$ & $12.9\%$ & $28.4\%$ \\
\bottomrule
\end{tabular}
\caption{Share of the total excess in Table~\ref{tab:excess},
  by how many observations the state has.}
\label{tab:excess-where}
\end{table}

Two facts follow.  The excess falls sharply with data per state:
\texttt{enwik8} and \texttt{enwik9} differ by a factor of nine in
observations per state and by a factor of three in excess, at a fixed
vocabulary and nearly fixed state count.  And the excess does not sit
in the near-empty states.  A state observed once costs $16.614$ bits
for its single symbol on every file, which is $\log_2 100{,}277$
exactly --- the model paying the uniform price over the whole alphabet
when it knows nothing.  (That figure is per symbol rather than per
character, because it is a statement about the alphabet; everything
else in this section is per character.)  Such states hold too few
tokens to matter.  Two thirds of the excess sits
in states with between one hundred and ten thousand observations
(Table~\ref{tab:excess-where}), where estimation would seem
comfortable.  A state with a thousand observations spread over a
hundred thousand symbols still pays, in effect, to say which symbols
occur, and it pays separately from every other state, although the
marginal distribution already answers most of the question.  What is
missing is a way for a state to borrow the marginal by a graded amount
while keeping its own evidence.  The family above cannot express that,
because the only alternative it offers to a state's own counts is to
discard them.

\section{Memory of order two}
\label{sec:ordertwo}

Each additional symbol of context multiplies the number of states
while the data stay fixed, and Table~\ref{tab:excess-where} prices the
consequence: a state observed twenty times carries about $2.2$ bits
per character of excess, against the $0.378$ paid on average, and
coarsening cannot recover the difference (Table~\ref{tab:coarsen}).
Whether a second symbol of context pays is therefore a question about
data per state, and it is answered by running it.

Keep full resolution on the previous token and add graded resolution
on the one before it: the state is $(x_{t-1},\, b_{M_2}(x_{t-2}))$,
where $b_{M_2}$ keeps a token if its vocabulary identifier is among
the $M_2$ smallest and buckets it otherwise, so that $M_2 = 0$
recovers the first-order model of Table~\ref{tab:firstorder} exactly.
The map must not depend on the data, which is why the retained set is
defined by identifier and not by frequency; a frequency threshold
would have to be transmitted.  Table~\ref{tab:order-two} gives the
result on all three files.

\begin{table}[ht]
\centering
\begin{tabular}{lrrrrr}
\toprule
file & $M_2$ & states & obs./state & bits/char & vs.\ order one \\
\midrule
\texttt{text8}
 &    0 &    35{,}767 &   543 & \textbf{1.8298} & --- \\
 &   16 &    59{,}664 &   326 & 1.8329 & $+0.0031$ \\
 &   64 &   179{,}826 &   108 & 1.8581 & $+0.0282$ \\
\midrule
\texttt{enwik8}
 &    0 &    71{,}160 &   363 & \textbf{2.1135} & --- \\
 &   16 &   146{,}050 &   177 & 2.1240 & $+0.0105$ \\
 &   64 &   243{,}396 &   106 & 2.1284 & $+0.0149$ \\
 &  256 &   427{,}694 &    60 & 2.1519 & $+0.0383$ \\
\midrule
\texttt{enwik9}
 &    0 &    81{,}275 & 3{,}367 & 1.8371 & --- \\
 &  256 & 1{,}063{,}906 &  257 & 1.7972 & $-0.0399$ \\
 & 1{,}024 & 2{,}753{,}260 &   99 & \textbf{1.7915} & $-0.0456$ \\
 & 2{,}048 & 3{,}719{,}340 &    74 & 1.7931 & $-0.0440$ \\
\bottomrule
\end{tabular}
\caption{Order two on the subword alphabet at the full vocabulary
  ($V = 100{,}277$, $M_1 = V$), against the first-order model
  ($M_2 = 0$) of Table~\ref{tab:firstorder}.  Observations per state
  are counted over states actually observed.  The sign of the answer
  is set by the corpus, not by the construction: on \texttt{text8} and
  \texttt{enwik8} every allocation of two-back resolution costs more
  to learn than it gains, and the family posterior puts mass $1.0$ on
  the first-order member; on \texttt{enwik9}, ten times the data, the
  same grid pays $0.046$ bits per character, with an optimum in the
  interior at $M_2 = 1{,}024$.  The two files that lose are exactly the
  two that begin with a few hundred observations per state; the one
  that wins begins with a few thousand.}
\label{tab:order-two}
\end{table}

\paragraph{Where the crossing is.}  Two files disagreeing is weak
evidence, because they differ in more than size.  Running prefixes of
\texttt{enwik9} varies size alone, and brackets the crossing on a
single source.  At $50$ million tokens the first-order model wins,
$7.7075$ against $7.7330$ bits per token at $M_2 = 64$; at $100$
million the order is reversed, $7.4277$ against $7.4182$; over the
whole file order two wins by $0.1666$.  In observations per first-order
state those three points are $669$, $1{,}283$ and $3{,}367$, so the
crossing lies between roughly seven hundred and thirteen hundred
observations per state.  (Prefixes are reported per token: bits per
character needs the whole file's byte count and the vocabulary charge,
neither of which applies to a prefix, so dividing a prefix's bits by
the file's bytes would give a number that is neither a rate nor
comparable between prefixes.)  Read together with
Table~\ref{tab:excess-where}, this is the same statement twice.  The
second symbol of context does not pay or fail to pay in itself; it pays
once the states it creates are large enough to learn, and the threshold
is a property of the estimator, measurable, and roughly a thousand
observations here.

The finding of these first three sections constrains everything that
follows: the cost of memory is set by what a single state must pay to
learn its own distribution, not by how many states there are.
Table~\ref{tab:firstorder} is the reference point for the
constructions of the next sections.

\section{Adaptive depth: the context-tree family}
\label{sec:ctree}

This section reports one construction in one measure.  The
construction is the hierarchy route of Appendix~\ref{app:routes}:
condition each position on its last $k$ symbols, with $k$ chosen per
context by an exact mixture over all prunings of the context tree up
to depth $D$, the layered estimator at every node, and the whole
codelength in closed form.  The measure is bits per character of the
original file, fully accounted: the tree's codelength over the LLM
token stream, plus the few boundary tokens the tree cannot
condition, plus the vocabulary charge of Section~\ref{sec:sources}
($0.0621$ bits per character at $10^{8}$ bytes).  At depth one the
construction \emph{is} the first-order model of
Section~\ref{sec:firstorder}, and its rows below reproduce
Table~\ref{tab:firstorder} to the last digit --- the accounting is
closed.

\paragraph{No chunking anywhere.}  Two different things could be
approximated here, and neither is.  The \emph{encoder} is fully
sequential with no staleness: at every position it predicts from
node tables that reflect every previously decoded token, updated
after every single token --- there are no checkpoints, no blocks, no
frozen tables.  The \emph{evaluation} of that encoder's codelength
is nevertheless exact and needs no walk over the file: because each
node's estimator is exchangeable, the product of its token-by-token
predictions telescopes to a closed form at the node's final counts,
and the mixture over prunings is summed exactly by the tree
recursion (Appendices~\ref{app:factorisation}
and~\ref{app:routes}).  The final counts are a device for
\emph{computing the length} of the sequential code, never an input
to any of its predictions --- the same property already used for
the memoryless and first-order models, extended here to the whole
tree family.  Every figure in Table~\ref{tab:ctw} is therefore the
exact codelength of the exact code: unlike the checkpointed
constructions of Section~\ref{sec:pairwise}, nothing in this
section pays, or approximates away, a staleness cost.

\begin{table}[ht]
\centering
\begin{tabular}{llrrr}
\toprule
file & scheme & bits/char & vs.\ first order & best published \\
\midrule
\texttt{text8}  & first order ($D{=}1$) & $1.8298$ & ---       & $1.398$ \\
\texttt{text8}  & tree, $D=2$           & $1.8185$ & $-0.0113$ & \\
\texttt{text8}  & tree, $D=3$           & $1.8182$ & $-0.0116$ & \\
\addlinespace
\texttt{enwik8} & first order ($D{=}1$) & $2.1135$ & ---       & $1.170$ \\
\texttt{enwik8} & tree, $D=2$           & $2.0676$ & $-0.0459$ & \\
\texttt{enwik8} & tree, $D=3$           & $2.0606$ & $-0.0529$ & \\
\addlinespace
\texttt{enwik9} & first order ($D{=}1$) & $1.8371$ & ---       & $0.858$ \\
\texttt{enwik9} & tree, $D=2$           & $1.6840$ & $-0.1531$ & \\
\texttt{enwik9} & tree, $D=3$           & $1.6598$ & $-0.1773$ & \\
\bottomrule
\end{tabular}
\caption{The context-tree family at full fidelity: honest bits per
  character of each file, vocabulary charge included, LLM
  tokenization throughout.  ``vs.\ first order'' is the change
  against the depth-one member, the first-order model of
  Table~\ref{tab:firstorder}.  Every complete fixed depth beyond one
  loses badly (depth two alone: $2.3072$ on \texttt{text8},
  $2.5257$ on \texttt{enwik8}); all of the gain is the per-context
  choice of depth.}
\label{tab:ctw}
\end{table}

Three remarks.  First, adaptive depth makes real progress on the
scoreboard: $0.012$ bits per character on \texttt{text8}, $0.053$
on \texttt{enwik8}, and $0.177$ on \texttt{enwik9} below the
first-order numbers --- the best honest figures in this paper so
far, and the clearest demonstration yet that what memory the data
can afford is set by observations per context: tenfold data turns
depth two from a $0.05$ improvement into a $0.15$ one, with
$2.75$ million depth-two contexts kept.  Second, the file and the
data volume decide how much depth is worth: on \texttt{text8} depth
three adds almost nothing to depth two, on \texttt{enwik8}'s markup
it adds $0.007$, and on \texttt{enwik9} --- the same markup with
tenfold data --- $0.024$, with $2.8$ million depth-three contexts
kept.  The \texttt{enwik9} curve has not saturated at depth three;
depth four is the next measurement there.  Third,
the remaining distance to the published targets ($0.42$ and $0.89$
bits per character) is far larger than what deeper trees can close;
the routes with longer reach (Appendix~\ref{app:routes}) remain the
program.  An earlier version of this section reported the same
family on a word-level tokenization, where adaptive depth showed no
gain at full fidelity; those runs are superseded --- on subword
tokens, context depth carries information that word tokens absorb
into the symbol itself.

\section{Memory of order two from pairwise statistics}
\label{sec:pairwise}

Section~\ref{sec:ordertwo} located the failure of order-two states:
not the ceiling but the learning cost of one successor distribution
per context pair, $V^3$-scale where the data supports far less.
Section~\ref{sec:ctree} recovered part of the order-two value by
letting each context choose its depth.  This section attacks the
same target from the other side: how much of the order-two ceiling
can be reached while \emph{learning} only pairwise objects --- the
three tables $(x_t, x_{t-1})$, $(x_t, x_{t-2})$ and
$(x_{t-1}, x_{t-2})$, roughly $3V^2$ numbers against $V^3$ --- and
it is the first measured step of the constraint-model route of
Appendix~\ref{app:routes}.  The measured campaign at the standard of
this paper is being run as this draft stands; this section fixes the
experiment's design, which is itself the product of an exploratory
round, and the accounting its table will use.

\paragraph{The predictors.}  Every table any rule consumes is
estimated by the layered predictor --- the only estimator in this
paper.  The pair tables do not determine the conditional
$p(x_t \mid x_{t-1}, x_{t-2})$; the combining rule is the
experiment:

\begin{itemize}
\item \emph{Mixture}: $w_0\, m(y) + w_1\, \hat p(y \mid x_{t-1}) +
  w_2\, \hat p(y \mid x_{t-2})$ over a small grid of weights, $m$
  the unigram table.
\item \emph{Product}: $m(y) \prod_{\delta \in \{1,2\}}
  \bigl[\hat p(y \mid x_{t-\delta})/m(y)\bigr]^{\beta_\delta}$,
  normalized.  The point $\beta_1 = \beta_2 = 1$ is the exact
  conditional of the maximum-entropy joint consistent with the two
  target--lag pairs (Section~\ref{sec:calibration}); tempered
  exponents discount the double counting of redundant lags.
\item \emph{Calibrated product}: the conditional of the
  maximum-entropy joint consistent with all \emph{three} pair
  tables; redundancy between the lags is then priced automatically
  rather than by a tempering grid
  (Section~\ref{sec:calibration}).
\item \emph{References}: the first-order model (the scoreboard
  holder to beat), order-two counts with backoff, and the
  \emph{exact} layered order-two code --- a mixture, so its
  sequential code telescopes to one evaluation per context pair at
  the final counts, with no checkpoint staleness at all.
\end{itemize}

\paragraph{The protocol.}  None of the combined rules telescopes, so
they are coded sequentially with checkpointed tables: the stream is
cut into $C$ blocks, every table and calibration is computed from
the data strictly before the block, the block is coded, the tables
refresh.  Every member pays the same staleness, so the internal
comparison is fair; the exact order-two reference pays none, and the
gap between its exact and checkpointed forms measures whether $C$ is
large enough.  The checkpoints are placed \emph{geometrically} ---
small blocks early, growing by a fixed ratio --- not equally.  This
was decided by measurement against a scheme whose exact sequential
value is known in closed form (the first-order model): at $C = 32$,
equal blocks cost $0.170$ bits per token of staleness on the
\texttt{text8} stream at $V=1024$ where geometric blocks cost
$0.032$, and at $C = 8$ the ratio is nine to one.  The reason is
the learning curve: the tables change fastest over the first
tokens, so equal spacing wastes its refreshes where nothing moves
and starves the stretch where everything does.  With geometric
placement at $C = 32$ the residual staleness is a few hundredths of
a bit per token --- known, small, and stated next to any number it
affects.  Two further consequences of the exploratory round: the
share-nothing layered reference is damaged far more than smooth
members by any checkpointing ($0.198$ bits per token at $C = 32$
equal), so it appears only in its exact, telescoped form; and
members are reported with the checkpoint schedule they were
computed under, always.

\paragraph{The accounting.}  As everywhere now: full LLM vocabulary,
complete codes, bits per character of the original file with the
vocabulary charge included.  At the full vocabulary the pair tables
are held sparsely --- each row a shared background plus corrections
at observed symbols --- and mixtures and products cost only the
observed entries, since a product of such rows lives on the union of
their supports.  The one open feasibility question is the calibrated
member: its fitted factors densify under the multiplicative fit, and
if no sparse form is found it will be reported at a moderate
alphabet as a \emph{complete} code --- the catch-all symbol's
contents priced through a memoryless side channel --- so that even
that number lives on the one scoreboard.

\paragraph{What the exploratory round fixed.}  A reduced-alphabet
round (the mechanism arena; its raw figures are not file-comparable
and are not reported here) made three design decisions.  Combining
rules must consume layered tables: with crude count-smoothed tables
the ranking of rules inverts, and the calibrated construction ---
the principled one --- was the member most damaged by poor table
estimates and most improved by good ones; under layered tables it
overtook the backoff reference outright.  Checkpointing at $C = 32$
is adequate for the smooth members and ruinous for share-nothing
layered ones (the $0.198$ above).  And the calibration fit needs its
margins made consistent to high precision before fitting, or it
stalls short of its optimum --- the fitter, not the model, sets the
floor.  The measured table --- the analogue of Table~\ref{tab:ctw}
for this construction, three files, honest bits per character ---
will replace this paragraph when the campaign completes.

\section{The calibrated pairwise predictor}
\label{sec:calibration}

Let $X_0$ be the target symbol and $X_1, X_2$ the two context
symbols.  Suppose the pair distributions $p(X_0, X_1)$ and
$p(X_0, X_2)$ are known.  These do not determine the conditional
$p(x_0 \mid x_1, x_2)$: many joints share the given pairs and they
disagree about the conditional.  The standard resolution is maximum
entropy: among all consistent joints, take the one that adds nothing
else.  For log-loss this choice is also minimax over the consistent
set.

\paragraph{Two pairs.}  Both constraints involve $X_0$, so the
maximum-entropy joint factorizes as
$p(x_0)\,p(x_1 \mid x_0)\,p(x_2 \mid x_0)$ and its conditional is
closed-form:
\begin{equation}
\label{eq:star}
  p^\ast(x_0 \mid x_1, x_2)
  \;\propto\;
  p(x_0)\,
  \frac{p(x_0 \mid x_1)}{p(x_0)}\,
  \frac{p(x_0 \mid x_2)}{p(x_0)}.
\end{equation}
This is the product rule of Section~\ref{sec:pairwise} at
$\beta_1 = \beta_2 = 1$.  It is exact if and only if $X_1$ and $X_2$
are conditionally independent \emph{given} $X_0$; marginal
independence of the context symbols is neither necessary nor
sufficient.  When the context symbols carry overlapping information
about the target, \eqref{eq:star} multiplies the same evidence twice.

\paragraph{Three pairs.}  Suppose additionally $p(X_1, X_2)$ is
known.  The maximum-entropy joint is now a pairwise field
\begin{equation*}
  p^\ast(x_0, x_1, x_2)
  \;\propto\;
  \psi_{01}(x_0, x_1)\,\psi_{02}(x_0, x_2)\,\psi_{12}(x_1, x_2),
\end{equation*}
with the potentials determined by the requirement that all three pair
margins of $p^\ast$ equal the given tables.  Two facts follow.
First, since $x_1$ and $x_2$ are observed at prediction time, the
factor $\psi_{12}$ is a constant and cancels: the conditional still
factorizes over the lags,
\begin{equation}
\label{eq:calibrated}
  p^\ast(x_0 \mid x_1, x_2)
  \;\propto\;
  \psi_{01}(x_0, x_1)\,\psi_{02}(x_0, x_2).
\end{equation}
Second, the potentials are no longer the pairwise ratios of
\eqref{eq:star}: they solve a coupled system in which the third
table participates, and the constraint has a direct reading.  The
margin conditions state that the product \eqref{eq:calibrated},
averaged over how the two context symbols actually co-occur (that
is, over $p(x_2 \mid x_1)$), reproduces each pairwise conditional;
in \eqref{eq:star} the same average is taken as if the context
symbols co-occurred only through $X_0$.  Redundancy between the lags
is thereby discounted in the fitted factors themselves.

The extreme case is instructive.  If $X_2$ is a copy of $X_1$, the
star rule \eqref{eq:star} squares the evidence, while the calibrated
conditional \eqref{eq:calibrated} collapses to
$p(x_0 \mid x_1)$ exactly: the duplicate lag is priced at zero.  If
instead the lags are conditionally independent given the target, the
$\psi_{12}$ constraint is inactive and \eqref{eq:calibrated} reduces
to \eqref{eq:star}.  Between the extremes it interpolates, which is
the role the tempering exponents $\beta_\delta < 1$ play by hand,
with one scalar where the calibration uses a full table.

The sparse factorization, intersection representation, solver, checkpoint
workflow, and numerical validation are described in
Appendix~\ref{app:calibration-implementation}.

Controlled full-file measurements for this model are in progress.  Their
implementation and complete-code accounting are specified in
Appendix~\ref{app:calibration-implementation}.

\appendix

\section{Implementation of the calibrated pairwise predictor}
\label{app:calibration-implementation}

Fix the observed-support context law and write
\begin{equation}
  q(y,a,b)=\widehat p_{AB}(a,b)q(y\mid a,b),
  \qquad
  q(y\mid a,b)\propto
  \exp\{\ell_0(y)+c_1(y,a)+c_2(y,b)\}.
  \label{eq:calibration-implementation}
\end{equation}
Normalization of the conditional analytically supplies the context-only
potential, so the implementation never stores a $V^3$ joint.  The two
correction families are chosen by the regularized convex calibration problem
described in Section~\ref{sec:calibration}; active $YA$ and $YB$ cells are
soft finite-sample targets, while the unigram and observed-support $AB$ laws
form the hard reference measure.

\paragraph{Observed-support approximation.}
Although the layered estimator supplies pair probabilities without a dense
table, exact calibration against a positive $\widehat p_{AB}$ would still
average over all $V^2$ context pairs.  Calibration is therefore restricted
to context pairs observed in the decoded prefix, and $\widehat p_{AB}$ is
renormalized on that support.  Separately estimated target--lag laws need not
be compatible with this context law at finite sample size; their active cells
are consequently soft constraints rather than being projected to compatible
tables.  For an unseen context pair, prediction uses the two-pair rule
\eqref{eq:star}.  The relaxation strength and fallback are fixed parts of the
coding rule.

For sparse evaluation define
\[
 r_1(y,a)=e^{c_1(y,a)}-1,
 \qquad r_2(y,b)=e^{c_2(y,b)}-1,
\]
and
\[
 S_{1a}=\sum_y p_0(y)r_1(y,a),
 \qquad S_{2b}=\sum_y p_0(y)r_2(y,b).
\]
Then the conditional normalizer on an observed context edge is
\[
 Z_{ab}=1+S_{1a}+S_{2b}
       +\sum_y p_0(y)r_1(y,a)r_2(y,b).
\]
Only the intersection of the active correction rows is needed for the final
term.  A compiled, memory-mapped intersection graph stores these incidences
once across checkpoints.  Margin and stochastic-gradient evaluations traverse
the relevant graph blocks directly; neither dense $V\times V$ pair tables nor
a triple table are constructed.

At checkpoint $k$, independent tasks construct the unigram law and the two
layered target--lag laws from the decoded prefix, assemble the calibration
problem, publish the required graph delta, fit the factors, and score the next
interval.  Fits currently use independent cold starts selected from a fixed
portfolio containing the unigram, each one-pair solution, their midpoint, and
their product.  The stochastic dual solver publishes its best finite
normalized iterate; an exact stationarity calculation is retained as a
diagnostic rather than used to inspect the block being encoded.  Checkpoints
and pair-estimation tasks are scheduled concurrently subject to their causal
dependencies.

Correctness tests compare the sparse conditional and its margins with literal
dense enumeration at small alphabets, exercise support growth and coincident
lags, and verify normalization on saved large-alphabet states.  End-to-end
corpus results are accepted only from a fresh pipeline that rebuilds counts,
pair estimates, calibration problems, factors, and interval scores under one
orientation convention.

\section{The bias of the empirical conditional entropy}
\label{app:mm}

The ``best possible'' column of Table~\ref{tab:firstorder} and the
target in Table~\ref{tab:excess} are computed from the file's own
counts, and such an estimate is systematically too small.  This
appendix states the correction used and why it is needed, since the
size of the bias differs between the files and a comparison across
them is otherwise not a comparison.

Let a state $s$ occur $n_s$ times with successor counts
$(c_{s1},\dots,c_{sk_s})$ over the $k_s$ symbols that actually follow
it, and write $\hat p_{sj} = c_{sj}/n_s$.  The empirical conditional
entropy is
\[
  \hat H(X \mid \text{prev})
  \;=\; \sum_s \frac{n_s}{n} \, \hat H_s ,
  \qquad
  \hat H_s = -\sum_{j} \hat p_{sj} \log_2 \hat p_{sj} .
\]
Each $\hat H_s$ is the entropy of the observed proportions rather than
of the true ones.  Sampling noise always makes a distribution look
less uniform than it is --- some symbols happen not to occur, others
occur more often than their probability warrants --- and entropy is
concave, so $\mathbb{E}\,\hat H_s < H_s$.  Expanding $\hat H_s$ about
the true probabilities and taking expectations gives the classical
first-order term
\[
  \mathbb{E}\,\hat H_s \;=\; H_s - \frac{k_s - 1}{2 n_s} + O(n_s^{-2})
  \quad \text{nats},
\]
so adding $(k_s-1)/2n_s$ nats to each state removes the leading bias.
This is the Miller--Madow correction.  Weighting by $n_s/n$ and
converting to bits, the correction to the whole table is
\[
  \frac{1}{2 n \ln 2} \sum_s (k_s - 1)
  \;=\; \frac{P - S}{2 n \ln 2},
\]
with $P$ the number of distinct (state, successor) pairs and $S$ the
number of states.  It depends on the data only through those two
counts.

Three properties matter for how the corrected figures are read.  The
correction grows with the number of distinct pairs, so a model with
more states is flattered more by the uncorrected estimate; this is
exactly the comparison made between \texttt{enwik8} and
\texttt{enwik9}, and leaving it out would exaggerate the difference
between them.  It is $0.026$, $0.027$ and $0.011$ bits per character on
\texttt{text8}, \texttt{enwik8} and \texttt{enwik9}, so it is small
against the excesses of $0.336$, $0.378$ and $0.131$ that it adjusts
--- the gap those tables report is not an artefact of the estimator.
And it is a first-order term only: it is accurate when $n_s$ is large
relative to $k_s$ and understates the bias when it is not, so for
states seen once or twice the corrected target is still optimistic.
Those states hold a negligible share of the tokens
(Table~\ref{tab:excess-where}), so this does not affect the totals,
but it means the corrected entropy should be read as a lower bound on
what a first-order model could achieve, not as an attainable target.

Nothing in the model or in any codelength depends on this correction.
It is applied only to the empirical quantity the codelengths are
compared against.

\section{Why the product over states is a code}
\label{app:factorisation}

Every codelength reported here for a model with memory is computed as a
product over states,
\[
  Q^{\sigma}(x^n) \;=\; \prod_{s} q\bigl(\lambda_s\bigr),
\]
with $\lambda_s$ the final count profile of the successors emitted from
state $s$, and $q$ the depth-averaged layered mixture.  Nothing is
recomputed as the file is read.  That deserves an argument, because it
appears to use counts the decoder does not have, and because the
product does not in general commute with the averaging inside $q$.

\paragraph{What is required of $q$.}  Let $q$ assign a probability to
each finite sequence over the alphabet, and require two properties.  It
is \emph{exchangeable}: $q$ depends on a sequence only through its count
profile.  And it is \emph{consistent}:
$\sum_{y} q(\lambda \oplus y) = q(\lambda)$ for every profile $\lambda$,
where $\lambda \oplus y$ is $\lambda$ with the count of $y$ raised by
one.  Both hold for $q^{(L)}$ at each depth, being the law of an
exchangeable sequence, and both are preserved by convex combination, so
they hold for $q_{\mathrm{avg}} = L_{\max}^{-1}\sum_L q^{(L)}$.

\paragraph{The state map.}  Let $\sigma$ send a past $x^{t-1}$ to a
state $\sigma(x^{t-1})$, using only symbols the decoder has already
recovered.  Write $s_t = \sigma(x^{t-1})$, and let $\lambda_s^{(t)}$ be
the profile of the successors emitted from state $s$ strictly before
time $t$.  Both are functions of the past alone.

\paragraph{Claim.}  Define the one-step predictor
\[
  P(y \mid x^{t-1}) \;=\;
  \frac{q\bigl(\lambda^{(t)}_{s_t} \oplus y\bigr)}
       {q\bigl(\lambda^{(t)}_{s_t}\bigr)} .
\]
Then $P$ is a probability distribution at every step, and
$\prod_{t} P(x_t \mid x^{t-1}) = Q^{\sigma}(x^n)$.

\paragraph{Proof.}  That $P$ sums to one over $y$ is consistency, applied
at the profile $\lambda^{(t)}_{s_t}$, which the decoder can form.  So
$\prod_t P(x_t \mid x^{t-1})$ is a probability assignment on sequences,
and $-\log_2$ of it is a codelength realisable symbol by symbol.  For
the identity, group the product by state.  Fix $s$ and let
$t_1 < \dots < t_{n_s}$ be the times at which the state is $s$.  Before
$t_1$ that state has emitted nothing, so $\lambda^{(t_1)}_s$ is empty;
and $\lambda^{(t_{i+1})}_s = \lambda^{(t_i)}_s \oplus x_{t_i}$, because
between $t_i$ and $t_{i+1}$ no successor is emitted from $s$.  Hence
\[
  \prod_{i=1}^{n_s}
  \frac{q\bigl(\lambda^{(t_i)}_s \oplus x_{t_i}\bigr)}
       {q\bigl(\lambda^{(t_i)}_s\bigr)}
  \;=\; \frac{q(\lambda_s)}{q(\emptyset)} \;=\; q(\lambda_s),
\]
the intermediate factors cancelling in pairs.  Multiplying over states
gives the claim. \qed

Exchangeability is what makes the telescoped numerator depend only on
the profile: the order in which a state's successors arrived is
irrelevant, so a running profile per state is all the decoder carries.
The final counts appear in our arithmetic and never in the decoder's.

\paragraph{The averaging does not commute with the product.}  Here a
plausible-looking alternative is wrong.  The two expressions
\[
  \prod_s \Bigl( \tfrac{1}{L_{\max}} \sum_L q^{(L)}(\lambda_s) \Bigr)
  \qquad\text{and}\qquad
  \tfrac{1}{L_{\max}} \sum_L \prod_s q^{(L)}(\lambda_s)
\]
are not equal.  Both are legitimate codes and they are different models:
the first draws a depth independently for each state, the second draws
one depth and applies it everywhere.  We compute the first, which is the
model of Section~\ref{sec:firstorder} --- the construction gives every
state its own weights and therefore its own depth.  The proof above
establishes the first, and nothing in it licenses the second.

The family average sits outside all of this.  A convex combination
$\sum_\sigma w_\sigma Q^{\sigma}$ of probability assignments on the same
sequence space is again one, and its one-step predictor is the mixture
of the members' predictors weighted by their \emph{posteriors} given the
past.  Its receiver keeps one bookkeeping per member --- that member's
own partition and its own running profiles --- predicts with each, and
mixes.  The posterior weights move at every step, yet the accumulated
cost is $-\log_2 \sum_\sigma w_\sigma Q^{\sigma}(x^n)$ with the
\emph{prior} weights on the final member joints, because each posterior
weight cancels against the member joint it is built from.

\paragraph{Which scheme the number measures.}  Every codelength in
this paper is obtained from one evaluation at the final counts.  It is
worth stating as plainly as possible what that number is, because more
than one sequential scheme can be described as ``averaging the
estimators'', and they do not agree.

Write $P^{(L)}(y \mid \lambda) = q^{(L)}(\lambda \oplus y)/q^{(L)}(\lambda)$
for the prediction the depth-$L$ model makes at a state whose running
profile is $\lambda$.  Two receivers suggest themselves:
\[
  \text{(a)}\;\; \sum_L w_L(\lambda)\, P^{(L)}(y \mid \lambda),
  \qquad
  \text{(b)}\;\; \frac{1}{L_{\max}} \sum_L P^{(L)}(y \mid \lambda),
\]
where $w_L(\lambda) = q^{(L)}(\lambda) / \sum_{L'} q^{(L')}(\lambda)$ is
the posterior over depth given what that state has seen so far.
Receiver (a) weights the depths by how well each has explained that
state's data; receiver (b) weights them equally at every step,
regardless.  Both sum to one over $y$, so both are codes.  They charge
different numbers for the same file.

Three statements, and the whole matter rests on them.

\emph{First: our number is receiver (a), not receiver (b).}  The
weighted sum in (a) is algebraically identical to
$q_{\mathrm{avg}}(\lambda \oplus y)/q_{\mathrm{avg}}(\lambda)$, a ratio
of one function at consecutive profiles, so the per-step costs cancel
in pairs and the total collapses to $-\log_2 \prod_s
q_{\mathrm{avg}}(\lambda_s)$ at the final profiles.  Receiver (b) is
\emph{not} such a ratio, nothing cancels, and its total must be
accumulated position by position.  It cannot be recovered from the
final counts at any price.

\emph{Second: the receiver never sees a final count.}  At time $t$ it
holds $\lambda$, the profile of what it has already decoded from the
current state, and forms both $q_{\mathrm{avg}}(\lambda)$ and
$q_{\mathrm{avg}}(\lambda \oplus y)$ from it.  The final profile is the
\emph{outcome} of its walk, not an input to it, and the closed form is
our shortcut for summing the walk --- verified on a $200{,}000$-token
prefix of \texttt{text8} over the subword alphabet, where the walk with
running counts and the closed form at the final counts agree to
$0.000$ bits.

\emph{Third: (a) is the better code, and that is why we want it.}  For
components $c$ with prior $\pi_c$,
\[
  -\log_2 \sum_{c} \pi_c Q^{(c)}(x)
  \;\le\; -\log_2 Q^{(c)}(x) + \log_2 (1/\pi_c)
  \qquad\text{for every } c,
\]
because a sum is at least any one of its terms.  With a uniform prior
this is the guarantee quoted throughout: never more than $\log_2$ of
the component count worse than the best component, on any file, chosen
with hindsight.  Receiver (b) admits no such bound and can be worse
than every component.  Measured on a byte stream, (b) costs
$0.049$ bits per character more than (a).

So the cheapest computation and the strongest scheme are the same one.
That is not quite a coincidence: both properties follow from the object
being a genuine mixture of joint distributions rather than a blend of
predictors.  The telescoping is that mixture's chain rule; the bound is
the fact that a sum exceeds its terms.  What is fortunate is only that
one had no obligation to be cheap.

\paragraph{How well consistency actually holds.}  The argument needs
$\sum_y q(\lambda \oplus y) = q(\lambda)$ exactly, and $q$ is a
numerical evaluation of an integral, so it holds only to the estimator's
accuracy.  That shortfall was measured rather than assumed.  Summing the
assigned probability over every sequence of a given length, at alphabets
small enough to enumerate exhaustively, the total is $0.987$ at $d=2$,
$0.992$ at $d=3$ and $0.994$ at $d=4$, improving with both alphabet size
and sequence length.  At the alphabet used here, $d = 100{,}277$, the
one-step predictor sums to within $2 \cdot 10^{-3}$ of one for a state
with five or more observations, and within $3.6 \cdot 10^{-2}$ for a
state seen once.  Accumulated over the $71{,}161$ states of the
first-order model on \texttt{enwik8}, the total effect is of order
$10^{-5}$ bits per character, four orders below the last digit reported
anywhere here.  Where the assignment is sub-normalised the reported
codelength is an over-estimate, which is the safe direction; the
measured deviations carry both signs and do not accumulate
systematically.

\section{The table of $\phi$ values: what is stored, what is
computed, and how it was checked}
\label{app:store}

Every codelength in this paper ends in the same place.  The layer
recursion of~\cite{layered2026} turns the block probability of a
count profile into a sum of one-dimensional integrals, and the
integrands are built from a single special function,
$\phi_r^{(L)}(t)$, defined in Appendix~B of~\cite{layered2026}.  Its
arguments are a level $L$ (one of the depths of the layered mixture),
a count $r$ (how many times some symbol occurred), and a point $t$ on
the integration axis.  Nothing else enters: not the alphabet, not the
file.  So a value of $\ln \phi_r^{(L)}(t)$, once computed, is correct
for every run, forever.  Evaluating a codelength means requesting
many such values; this appendix says how those requests are served
and why the answers are trusted.

\paragraph{The four ways we can compute $\phi$.}
We have four independent methods, with different strengths.
(1)~At $L = 1$ there is an exact closed formula.
(2)~A \emph{contour integral} (numerical integration of an exact
integral representation) is accurate everywhere but slow.  Because it
shares no code and no approximation with the other methods, it is our
reference: whenever we say a value was checked, we mean checked
against the contour integral.
(3)~\emph{Series expansions} converge at the two ends of the $t$
axis.  Each series evaluation returns two numbers: the value, and a
computed bound on its own error.  If the bound is not tiny, the value
is thrown away and the contour integral is used instead.
(4)~A \emph{saddlepoint approximation} is nearly free to evaluate and
needs no storage.  Its accuracy improves with the level: its largest
error along the axis is about $7 \cdot 10^{-3}$ bits at $L = 2$ and
shrinks to about $3 \cdot 10^{-6}$ bits by $L \approx 35$,
essentially independent of $r$.  The largest error is the number that
matters, because a codelength adds up thousands of $\phi$ values from
all over the axis, and an error anywhere on the axis can land in that
sum.  So the approximation is unusable at shallow levels and
excellent at deep ones; we exploit exactly that below.

\paragraph{What is stored.}
For every level from $2$ to $53$ and a designed set of counts, we
store a \emph{column}: the values of $\ln \phi_r^{(L)}$ at a regular
grid of points along the axis (spacing $0.02$ in $\ln t$), computed
by the contour integral and the series with their error bounds
enforced.  A request for a point between grid points is answered by a
degree-seven polynomial through the neighbouring grid values.  The
set of counts is chosen to match how counts occur in files.  In any
file, most distinct symbols are seen only a few times, so small
counts are requested constantly: every integer count from $0$ to
$255$ has its own column, and requests there involve no approximation
across counts at all.  Large counts belong to a handful of very
frequent symbols, and their exact values vary from file to file, so
storing a column for every large integer would be enormous and
almost never reused.  Instead, above $255$ we store columns only at
\emph{anchor} counts --- rounded powers of $1.083$, that is, anchors
about $8\%$ apart, up to $r = 2 \cdot 10^8$ --- and a request for a
count between anchors is answered by a degree-eleven polynomial in
$\ln(r{+}1)$ through the twelve nearest anchor columns.  The whole
table occupies $2.8$ gigabytes.  It is read-only: nothing appends to
it, a request it cannot serve stops the run with an error naming the
missing column, and every column carries a checksum that is verified
each time it is read.

\paragraph{What is computed instead of stored.}
Level $1$ uses its closed formula.  Levels $54$ and above store
nothing: there the saddlepoint approximation (with the series at the
ends of the axis) answers directly, since by those depths its largest
error is far below anything a reported number can feel.  The cutoff
$54$ was chosen by measurement, not judgment: serving levels from
$46$ up by the approximation shifts the first-order \texttt{text8}
codelength by $2.4 \cdot 10^{-4}$ bits per token, enough to change
the last reported digit; extending the stored levels to $53$ costs
$0.4$ gigabytes and the digit holds.

\paragraph{How the whole arrangement was checked.}
The requirement is stated in the unit we actually report, bits of
codelength, not in the accuracy of individual $\phi$ values.  We
fixed $22$ test count profiles chosen to cover everything the
paper's experiments produce --- largest counts from $200$ to
$1.7 \cdot 10^7$, alphabets up to $300{,}000$, mixture depths up to
$53$ --- and required that the codelength computed through the table
agree, on every one of them, to within $10^{-4}$ bits with the
codelength computed from exact columns built fresh by the contour
integral.  Measured: the worst disagreement is $5.6 \cdot 10^{-6}$
bits, eighteen times inside the requirement.  As an end-to-end check,
the first-order codelength of \texttt{text8} computed both ways
differs by $2.2 \cdot 10^{-5}$ bits per token; no reported digit
moves.  The safety margin was also measured: removing every second
anchor makes exactly one of the $22$ profiles fail, so the anchor
density stands about a factor of two from the edge.  Two precautions
guard the reference itself: the exact columns used in these
comparisons are kept in a separate frozen table, each column verified
identical to a fresh contour evaluation; and inside the reference the
series is disabled, so a defect in the series could not confirm
itself --- a rule adopted after precisely such a defect (a series
whose error bound was too optimistic in one regime) was caught and
repaired.  With the table in place, the paper's complete set of model
runs re-executes in a few minutes on a laptop.

\section{Four routes to long memory, in the same calculus}
\label{app:routes}

Every construction in this paper is one move made repeatedly: enlarge
the model by a latent variable, put a prior on it, and integrate the
likelihood exactly --- over the simplex point $\theta$, over the depth
$L$, over the family member.  Nothing is fitted; the posterior does
whatever adaptation there is.  This appendix records four known
constructions that extend the same move to memory far beyond fixed
order.  They are candidates for the next stage of the program, not
implemented here.  Each is presented the same way.  First the
prediction: we stand at a position of the sequence and say exactly
what is predicted from what.  Then the codelength: the trick, where
one exists, that computes the total without walking the file.  Then
the cost, and the evidence.  The four routes enlarge the latent
space along different axes --- tree depth, the set of statistics the
model is asked to respect, a copy pointer, time --- and compose with
one another.

\paragraph{Route 1 --- Hierarchy: context-tree weighting.}

\emph{Prediction.}  We have a sequence $x_1, x_2, \dots, x_n$.  We
stand at position $t$: the symbols $x_1, \dots, x_{t-1}$ are decoded,
and $x_t$ is to be predicted.  Fix a maximal depth $D$ and read
backwards from position $t$.  This gives $D{+}1$ descriptions of
where we stand, one per depth.  At depth $0$ the description is
empty.  At depth $1$ it is the last symbol, $x_{t-1}$.  At depth $2$
it is the last two symbols, and so on, up to depth $D$.  Call the
depth-$k$ description the depth-$k$ \emph{context} of position $t$.

Every string that occurs as a context gets its own estimator --- for
us, the layered average of Section~\ref{sec:setup}.  Its data are
simple.  Each time the string appeared in the past, some symbol came
next; the estimator counts these successors.  At position $t$ it has
seen only the occurrences inside $x_1, \dots, x_{t-1}$, and it
predicts from those counts, like every estimator in this paper.

So at position $t$ we hold $D{+}1$ predictions of $x_t$, one per
depth, and they differ in a familiar way.  The depth-$0$ estimator
has seen $t{-}1$ symbols but ignores the context entirely.  The
depth-$D$ estimator fits the current situation exactly but may have
seen it only a handful of times.  Some depth is best, and the best
depth differs from context to context.  Forcing one global depth is
what Table~\ref{tab:coarsen} measured, and it was too rigid.

Context-tree weighting (CTW)~\cite{ctw95} therefore makes the depth a
latent variable, local to each context, and integrates it out.
Arrange all context strings as a tree: the empty string is the root,
and the children of a string extend it by one more symbol into the
past.  A \emph{pruning} is a cut of this tree; it prescribes, for
every branch, the depth at which conditioning stops.  Under a given
pruning, each position of the file is predicted by exactly one node:
the deepest of its $D{+}1$ contexts that the pruning kept.  Put the
prior $2^{-\Gamma(T)}$ on prunings $T$, where $\Gamma(T)$ counts the
nodes of $T$, and predict $x_t$ by the posterior average over all
prunings.  The posterior weight of $T$ is, as always, its prior times
the probability it has assigned to $x_1, \dots, x_{t-1}$.

There are exponentially many prunings, so averaging over all of them
sounds hopeless.  It is not, for a simple reason: as far as position
$t$ is concerned, a pruning matters only through the one context it
keeps there.  Two prunings that cut the path of position $t$ at the
same depth predict $x_t$ identically, however much they differ
elsewhere in the tree.  So the average over all prunings is really an
average over $D{+}1$ alternatives --- predict from the depth-$0$
context, or the depth-$1$ context, and so on --- each weighted by the
combined posterior of all the prunings that cut there.  This average
can be computed from the deepest context upward.  Let $p_k$ be the
prediction of the depth-$k$ context's estimator, and blend:
\begin{equation*}
  \tilde q_{D} = p_{D},
  \qquad
  \tilde q_{k}
  \;=\;
  \gamma_k\, p_{k} + (1-\gamma_k)\, \tilde q_{k+1},
  \quad k = D{-}1, \dots, 0.
\end{equation*}
The final blend $\tilde q_0$ is the prediction.  In words: start
from the deepest context's prediction and, moving toward the root,
mix in each shallower context's prediction with weight $\gamma_k$.
The weight $\gamma_k$ is not a constant to be tuned.  It is a
posterior, computed from the record of the depth-$k$ node: how well
that node's own predictions have done, on the positions it has seen,
against the predictions of its descendants.  Coding $x_t$ updates
the $D{+}1$ estimators on its path and nothing else.  Predictors of
the PPM family postulate a rule of exactly this shape --- fall back
from deep contexts to shallow ones --- as a heuristic.  Here nothing
is postulated: the rule with these particular weights \emph{is} the
exact average over prunings.

\emph{The codelength.}  The scheme is sequential, and nothing it uses
at position $t$ lies beyond $x_1, \dots, x_{t-1}$.  Its codelength is
the usual sum, $-\sum_t \log_2 \tilde q_0(x_t)$, and as everywhere in
the paper the sum has a closed form.  Two telescopes produce it.  The
first works along the file, one node at a time.  A node's estimator
is exchangeable, so the product of its step-by-step predictions
equals its block probability at the node's \emph{final} successor
counts --- one evaluation of the layer recursion, served by the
table of Appendix~\ref{app:store}, by the argument of
Appendix~\ref{app:factorisation}.
Write $P_s$ for that number at node $s$.  The second telescope works
across prunings.  Define, from the leaves up,
\begin{equation*}
  Q_s \;=\;
  \begin{cases}
    P_s, & |s| = D,\\[2pt]
    \tfrac12\, P_s \;+\; \tfrac12 \prod_{a} Q_{as}, & |s| < D,
  \end{cases}
\end{equation*}
the product over the symbols $a$ seen preceding $s$ in the file.
Then the total codelength is exactly $-\log_2 Q_{\varnothing}$: the
recursion sums prior times likelihood over every pruning in one sweep
of the tree.  The final counts enter only through these two
evaluations.  They are a device for computing the length of the
sequential code, never an input to a prediction.

\emph{Fit with our scheme.}  Exact.  CTW does not care what the node
estimator is, only that it is a valid sequential assignment; ours is,
and both telescopes are the paper's own.  The construction also names
what Table~\ref{tab:coarsen} was missing.  The coarsening family
charged $\log_2$ of the member count for one global depth choice.
Here every context chooses its own effective depth, at a structural
price of about one bit per node actually used, and shallow nodes see
pooled counts --- graded borrowing from the parent, built into the
estimator.  At $D=2$ the mixture contains memoryless, first and
second order per context: the Bayesian completion of the backoff rule
that won the pairwise round of Section~\ref{sec:pairwise}.

\emph{Cost.}  Nodes are created only when visited: at most $D$ new
nodes per position, at most $O(Dn)$ ever.  Time is one layered block
evaluation per node --- profiles repeat across nodes, so the store
amortizes them --- plus a linear sweep; memory is linear in the node
count.  The vocabulary enters only through the per-node estimator,
which is already ours.

\emph{The alphabet is large.}  Our alphabet is the tokenizer
vocabulary, $d \approx 10^5$, and a $d$-ary tree sounds hopeless:
there are $d^2$ possible contexts at depth two alone.  But the tree
never materializes.  A node is created the first time its context
occurs in the file, so after $n$ symbols at most $Dn$ nodes exist,
whatever $d$ is --- and far fewer in practice; on \texttt{enwik8}
only $71{,}161$ of the $d$ possible depth-one contexts occur at all
(Table~\ref{tab:firstorder}).  What limits useful depth is data, not
the alphabet.  At larger depths most contexts occur once or twice,
their nodes hold one or two counts, and the posterior weights
$\gamma$ shut them off in favour of their parents, at the structural
price the recursion has already charged --- about one bit per node
used.  Storage can even be made linear in $n$ independently of $D$,
by the standard path compression: a chain of nodes each having a
single observed child is stored as one edge, which is how PPM* and
the sequence memoizer run at unbounded depth.  So $d$ never enters
the size of the tree.  It enters only the per-node estimator, and
that estimator is the one this paper already runs at this $d$.

\emph{Synchronization.}  Using this in a transmission scheme costs
nothing beyond the codelength.  The tree, its counts, and the
blending weights are deterministic functions of the decoded prefix,
so the decoder rebuilds, just before position $t$, exactly the
predictor the encoder used at position $t$ --- the situation of every
scheme in this paper.  What encoder and decoder must share is the
rule, fixed in advance: $D$, the priors, and, as with any arithmetic
coder, bit-exact arithmetic.  That is part of the definition of the
code, $O(1)$, not a per-file overhead; no side information
accompanies the data.  The same holds for the routes below: their
fitted factors, match sites and switch weights are likewise functions
of the decoded prefix alone.

\emph{Evidence.}  Published: CTW achieves the optimal redundancy for
the class of tree sources~\cite{ctw95}, one of the few universal
schemes with a sharp guarantee; its unbounded-depth Bayesian
relative, the hierarchical Pitman--Yor chain of the sequence
memoizer~\cite{memoizer}, matches PPM-class compressors on text with
linear resources.  In-house: markov2-with-backoff beat every
share-nothing order-2 member in the pairwise round, and the learning
shortfall of Table~\ref{tab:firstorder} ($0.400$ bits per character
on \texttt{enwik8}, the largest single measured loss in the paper) is
precisely the quantity that parent-sharing attacks.

\paragraph{Route 2 --- Constraint models: maximum entropy from
pairwise and triple statistics.}
This route is the direct generalization of the experiment of
Section~\ref{sec:pairwise}, which is already running; it is listed
here because the generalization, not the first step, is the prize.

\emph{Prediction.}  We stand at position $t$ with
$x_1, \dots, x_{t-1}$ decoded.  Choose a set of lags, for example
$1, 2, 4, 8, \dots$: the context is the symbols at those distances
behind $t$.  From the decoded prefix, estimate low-order statistics
--- for each chosen lag, a pair table saying how often symbol $a$ is
followed at that distance by symbol $b$; if desired, pair tables
\emph{between} context positions too, and a few selected triples.
These statistics do not determine the joint distribution of the
target and its context: many joints share them.  The principle of
Section~\ref{sec:pairwise} picks one --- the joint of maximum entropy
among all that match the statistics, the one that assumes nothing
beyond them.  Its form is classical: one factor per constraint, a
\emph{factor graph} whose variables are the target and the context
positions and whose factors sit on the constrained pairs and triples.
The factors are fitted to reproduce the tables by iterative
proportional fitting, the same convex procedure the calibrated member
already uses; nothing in the fit ever sees the codelength.  And the
prediction is cheap even when the graph has loops: every context
symbol is observed, so the distribution of the target is proportional
to the product of the factors touching the target --- one sweep over
the alphabet, no inference beyond that.  With only target--lag pairs
the joint factorizes and this reduces exactly to the product rule at
exponent one; adding lag--lag constraints is the principled cure for
the double counting that made that rule lose, because redundancy
among the lags is then priced into the factors
(Section~\ref{sec:calibration} proves the two-lag case).  Which
constraints to include is not tuned either: it is a latent choice ---
a family of constraint sets, mixed with a prior, exactly as the
current family of members is mixed.

\emph{The codelength.}  As in the running experiment: the tables, and
the factors fitted from them, are refreshed at checkpoints, and the
codelength is accumulated sequentially by the machinery of
Section~\ref{sec:pairwise}, with the same measured exact-versus-%
checkpointed control.

\emph{Fit with our scheme.}  This \emph{is} our scheme, extended.
The calibrated member is this route with the smallest interesting
constraint set; Section~\ref{sec:calibration} is its two-lag theory;
the tables come from the same smoothed count estimators as everything
else.  What the generalization adds is reach: more lags at additive
cost, lag--lag constraints, selected triples where synergy is
suspected.

\emph{Cost.}  The statistics grow additively --- one $V \times V$
table per constrained pair, against the $V^{\Delta}$ of a Markov
model over the same span --- and the per-position prediction is a
product over the constraints touching the target.  At the full
vocabulary the tables are sparse and the product lives on the union
of the observed supports, as noted in Section~\ref{sec:pairwise};
whether the \emph{fitted factors} can be kept sparse as well is the
open question of that section, and it is the real obstacle between
this route and the full vocabulary.

\emph{Evidence.}  Published: maximum-entropy language models built
from exactly such constraints --- including long-distance
``trigger'' pairs --- were state of the art in their day and reduced
perplexity substantially over fixed-order models~\cite{rosenfeld};
fitting is convex, so the construction is as solid as its
statistics.  In-house, measured: the calibrated member beats every
mixture in the pairwise round with no free parameter, and the round's
best members recover about two thirds of the order-two gain from
pairwise objects alone.  The known limit is also measured: pair
statistics at long distances are thin (the pooled-lag tail was worth
$0.044$ bits per token), and pairwise constraints cannot express
synergy --- two lags that are individually uninformative but jointly
decisive.  Triples address the second point; nothing in this route
addresses the first, which is retrieval's territory.

\paragraph{A Bayesian completion of the constraint route.}
There is a more ambitious version of Route~2 that places it exactly
in the calculus of Section~\ref{sec:setup}.  It is useful to state it
now because it separates an exact statistical identity from the
computational obstacle that any implementation must overcome.  Let
\[
  Z_t=(X_t,X_{t-1},\ldots,X_{t-k})
\]
be the window ending at position $t$, and let $T(Z_t)$ collect the
selected pair and triple indicators.  A natural-parameter vector
$\eta$ defines the exponential-family joint
\begin{equation}
\label{eq:future-exp-family}
  p_\eta(z)
  =\exp\!\left\{\langle\eta,T(z)\rangle-A(\eta)\right\},
  \qquad
  A(\eta)=\log\sum_z
       \exp\!\left\{\langle\eta,T(z)\rangle\right\}.
\end{equation}
Matching prescribed pairwise marginals is the mean-parameter equation
$\nabla A(\eta)=\mu$.  The calibrated predictor currently solves one
small instance of this inverse problem.  The proposed reversal is not
to solve it anew for every possible $\mu$, but to put a prior $\Pi$ on
a region of natural parameters and average the resulting predictors.

If, as a first mathematical model, the $N$ observed windows are
treated as independent, their sufficient statistic is
\[
  S_N=\sum_{t=1}^{N}T(Z_t),
\]
the collection of empirical low-order count tables, and the Bayesian
assignment is exactly
\begin{equation}
\label{eq:future-bayes-joint}
  Q_N(Z^N)
  =\int \exp\!\left\{
       \langle\eta,S_N\rangle-NA(\eta)
     \right\}\,\Pi(d\eta).
\end{equation}
Consequently the complete sequential log loss telescopes to
\begin{equation}
\label{eq:future-bayes-loss}
  -\sum_{t=1}^{N}\log_2 Q(Z_t\mid Z^{t-1})
  =-\log_2 Q_N(Z^N).
\end{equation}
This is the direct analogue of the profile formula for the layered
simplex prior: the data enter only through $S_N$, and a prior with
additional structural indices --- constraint set, graph, rank or
interaction range --- can be averaged in the same expression.  For a
countable family $m$ with weights $w_m$,
\begin{equation}
\label{eq:future-structure-mixture}
  Q_N^{\rm mix}=\sum_m w_m Q_N^{(m)},
  \qquad
  -\log_2 Q_N^{\rm mix}
  \leq \min_m\{-\log_2 Q_N^{(m)}-\log_2 w_m\}.
\end{equation}
Thus uncertainty about which pairwise constraints to retain costs
only their description length, rather than a tuned selection on the
file.

Equation~\eqref{eq:future-bayes-joint} is exact but not yet an
algorithm.  Every value of its integrand requires the global
normalizer $A(\eta)$, a sum over the alphabet of the entire window;
Bayesian averaging adds an integral over $\eta$ but does not remove
that sum.  The formal conjugate prior makes the limitation especially
plain.  With
\[
  \Pi_{\alpha,\nu}(d\eta)
  =C(\alpha,\nu)^{-1}
    \exp\!\left\{\langle\eta,\alpha\rangle-\nu A(\eta)\right\}d\eta,
\]
the marginal likelihood is the compact ratio
\begin{equation}
\label{eq:future-conjugate}
  Q_N(Z^N)
  =\frac{C(\alpha+S_N,\nu+N)}{C(\alpha,\nu)},
\end{equation}
but the conjugate normalizer $C$ is generally as difficult as the
partition function it contains.  Conjugacy relocates the computation;
it does not solve it.  This is the central feasibility test for the
fully Bayesian constraint model.

There are three ways around the global normalizer, each worth a
controlled experiment and each defining an exact probability assignment
once its local ingredients are computed.  First, predict conditionally
rather than construct the whole window joint:
\begin{equation}
\label{eq:future-conditional-family}
  p_\eta(x_t\mid x_{t-1:t-k})
  =\frac{\exp\{\langle\eta,T(x_t,x_{t-1:t-k})\rangle\}}
         {\sum_y\exp\{\langle\eta,T(y,x_{t-1:t-k})\rangle\}}.
\end{equation}
The global $A(\eta)$ cancels and only a sum over candidate next tokens
remains.  A prior over $\eta$ then gives the coherent sequence code
\begin{equation}
\label{eq:future-conditional-mixture}
  Q(x_{k+1:n}\mid x_{1:k})
  =\int\prod_{t=k+1}^{n}
       p_\eta(x_t\mid x_{t-1:t-k})\,\Pi(d\eta),
\end{equation}
whose posterior predictive is the posterior, not prior, average of
the fixed-$\eta$ predictors.  This is also the correct resolution of
the overlap between consecutive windows: treating them as independent
is a useful likelihood approximation, whereas
\eqref{eq:future-conditional-mixture} is a probability assignment to
the original stream.

Second, restrict the constraint graph to a tree.  For mutually
consistent node and edge marginals, the maximum-entropy joint is
\begin{equation}
\label{eq:future-tree-joint}
  p(z)=
  \frac{\prod_{(i,j)\in E}p_{ij}(z_i,z_j)}
       {\prod_i p_i(z_i)^{\deg(i)-1}},
\end{equation}
already normalized, with no partition function to fit.  The same
property extends to decomposable graphs by replacing edges with
cliques and dividing by their junction-tree separators.  Rooting the
tree turns it into local conditional tables; Dirichlet or layered
priors on those tables make the integrated codelength a product of
per-table mixture probabilities, computable from the pair counts by
the factorization of Appendix~\ref{app:factorisation}.  A mixture over
trees or bounded-width decomposable graphs then pays the structural
penalty in \eqref{eq:future-structure-mixture}.

Between this exact restriction and the unrestricted graph lies an
approximate route: belief propagation and the Bethe free energy.  For
a pairwise factorization
\[
  p_\eta(z)=Z(\eta)^{-1}
    \prod_i\psi_i(z_i;\eta)
    \prod_{(i,j)\in E}\psi_{ij}(z_i,z_j;\eta),
\]
the messages
\begin{equation}
\label{eq:future-bp-message}
  m_{i\to j}(z_j)
  \ \propto\
  \sum_{z_i}\psi_i(z_i)\psi_{ij}(z_i,z_j)
  \prod_{\ell\in\partial i\setminus j}m_{\ell\to i}(z_i)
\end{equation}
produce node and edge beliefs $b_i,b_{ij}$.  On a tree they give the
exact marginals and exact partition function.  On a graph with loops,
fixed points of the same recursion are stationary points of the Bethe
free energy~\cite{yedidia}, and give the approximation
\begin{equation}
\label{eq:future-bethe}
  A(\eta)=\log Z(\eta)
  \ \simeq\
  A_{\rm B}(\eta):=-\min_{b\in\mathcal L}F_{\rm B}(b;\eta),
\end{equation}
where $\mathcal L$ is the locally consistent marginal polytope.  This
replaces a sum over every configuration of the window by local message
updates.  It is exact on trees; with loops, convergence and accuracy
must be measured rather than assumed.

The approximation can enter at two levels.  In the inverse problem,
the Bethe beliefs approximate $\nabla A(\eta)$, so the potentials can
be updated until $b_{ij}$ matches the prescribed pair tables.  In the
Bayesian calculation, one may replace $A$ by $A_{\rm B}$ in
\eqref{eq:future-bayes-joint}, and then approximate the remaining
parameter integral by a finite set of particles, variational Bayes or
a Laplace expansion around its posterior mode.  Consecutive
checkpoints change the tables gradually, so both the messages and the
parameter fit can be warm-started.  Sparse factors, low-rank messages
and a \emph{tree plus a few correction edges} family provide natural
scaling controls at the full alphabet.

There is an accounting distinction important to this paper, but it
should not be confused with the computational obstacle.  Substituting
$A_{\rm B}$ in a batch marginal likelihood does not by itself prove
that the resulting number is the probability of a realizable code.
An actual arithmetic coder instead needs, at position $t$, only a
normalized distribution over the next token.  Message passing may
construct nonnegative scores $\widetilde s_t(y)$, after which encoder
and decoder both compute
\begin{equation}
\label{eq:future-bp-code}
  \widetilde Q_t(y\mid x^{t-1})
  =
  \frac{\widetilde s_t(y)}
       {\sum_{y'}\widetilde s_t(y')},
  \qquad
  \widetilde Q(x^n)=\prod_t
       \widetilde Q_t(x_t\mid x^{t-1}).
\end{equation}
Then the accumulated loss is the exact codelength of the implemented
approximate-inference predictor, even though the messages and beliefs
inside it are approximate.  The encoder and decoder must of course use
the same deterministic message schedule, finite-precision convention
and conversion to arithmetic-coder frequencies, but no fitted model
has to be transmitted when every update uses only the decoded prefix.

The denominator in \eqref{eq:future-bp-code} is merely one pass over
the target alphabet, $O(V)$ for vocabulary size $V$.  Even at
$V=10^5$ or $10^6$ that single operation need not be troublesome, and
with checkpointed factors it can be cached for repeated contexts or
computed from a shared background plus sparse corrections.  It is not
the global partition function in \eqref{eq:future-exp-family}, whose
sum ranges over all configurations of the $(k+1)$ variables and has
nominal size $V^{k+1}$, nor is it the high-dimensional parameter
integral in \eqref{eq:future-bayes-joint}.  Those are the computations
for which Bethe, particles and variational approximations are being
proposed.  The explicit $O(V)$ normalization serves chiefly to turn
their approximate scores into an actual code.

The resulting sequential length generally no longer telescopes to one
final profile evaluation, but it meets the same honest-code standard
as every measured route here.  Its total naive normalization cost is
$O(nV)$ if a new score vector is formed at every position, so whether
it matters is an implementation question: checkpointing, repeated
contexts, vectorization and the sparse-background representation
already used in Section~\ref{sec:pairwise} may reduce it sharply.
This makes loopy BP a particularly useful experiment: compare its
sequential codelength and measured computation against the exact tree
backbone while adding correction edges until the gain saturates or the
messages cease to be reliable.

Third, retain the present checkpointed calibrated conditional as a
finite or discretized family.  Partition functions or fitted factors
can be computed once per checkpoint and member, after which posterior
weights over the members update sequentially.  This does not integrate
over a continuous parameter region, but it is a complete Bayesian code,
provides a computable benchmark for the continuous construction, and
reuses the infrastructure already validated in
Section~\ref{sec:pairwise}.

These alternatives answer different questions.  The unrestricted
joint model is the faithful formulation when all prescribed pairwise
marginals must be reproduced, but its normalizer is the obstacle.  A
tree or decomposable model preserves exact marginal semantics while
making normalization local; Bethe interpolation adds loops at the
price of an approximation whose error can be measured against that
backbone.  The conditional model gives up a joint law for all context
variables, but it targets the quantity this paper actually codes,
$p(\text{current}\mid\text{past})$, and is therefore the most direct
next experiment.  The decisive comparison is not only their final
codelength: it is codelength together with the computation needed to
update one posterior prediction at the full alphabet.

\paragraph{Route 3 --- Retrieval: a mixture over copy pointers.}

\emph{Prediction.}  Again we stand at position $t$ with
$x_1, \dots, x_{t-1}$ decoded.  Text repeats itself.  The stretch we
are in --- the last tens or hundreds of symbols --- has often
occurred before, verbatim, and when it has, the file frequently
continues as it did there.  No Markov estimator of any order can say
this; the record-holding compressors give it a mechanism of its own
(the ``match model'' of the PAQ family~\cite{paq}, the unbounded
contexts of PPM*~\cite{ppmstar}).  The Bayesian rendering, worked out
by Steinr\"ucken~\cite{steinruecken}, is a latent pointer.

Call a position $m < t$ a \emph{match site} if the symbols before $m$
agree with the symbols before $t$ over some length.  If the file is
currently inside a repeated stretch, the symbol that followed $m$ is
an excellent guess for $x_t$.  Let the hidden variable $a_t$ say
whether we are copying: its value is either $\mathsf{bg}$ --- the
background model predicts --- or a match site $m$ --- the file is
copying the continuation of $m$.  Given $a_t = m$, the prediction
puts probability $1-\varepsilon$ on the symbol that followed $m$ and
spreads $\varepsilon$ over the background.  Given
$a_t = \mathsf{bg}$, it is the background's prediction.  The pointer
evolves under a prior: it persists with high probability, so a long
copy, once believed, costs almost nothing per symbol; sometimes it
breaks back to $\mathsf{bg}$; sometimes it jumps to a newly available
site.  The decoder keeps one weight per candidate value --- the
posterior, prior dynamics times predictive record --- updates the
weights each step by the forward recursion (multiply by the
probability the candidate gave the symbol just decoded, apply the
transitions, renormalize), and predicts by the weighted average.  The
escape rate $\varepsilon$ is not tuned; it is integrated under a Beta
prior from its own successes and failures.  Everything used at
position $t$ is computable from the decoded prefix.

\emph{The codelength.}  Summing over pointer paths makes this one
valid assignment $Q(x^n)$.  But here no telescope exists.  The
posterior over pointers depends on the whole path, exchangeability is
lost, and the codelength must be accumulated by walking the file ---
precisely the checkpointed sequential machinery built for
Section~\ref{sec:pairwise}, whose exact-versus-checkpointed gap is
already a measured quantity there.  Match sites come from a suffix
automaton over the decoded prefix, which yields the current longest
matches in amortized constant time; keeping the $K$ longest suffices
(the PPM* observation~\cite{ppmstar}).

\emph{Fit with our scheme.}  The background \emph{is} the existing
model, unchanged; retrieval sits on top as a mixture at every step
and reduces to the current code when the copy weight is zero.  No
weight in it is fitted --- persistence, breaks and escapes are all
priors updated by counts.

\emph{Cost.}  $O(K)$ weight updates per position, $K$ small; the
suffix automaton is linear in $n$; and nothing depends on the
vocabulary size, so this route runs at the full vocabulary from the
start.  Memory is linear in $n$.

\emph{Evidence.}  Published: the long memory of every record holder
on \texttt{enwik8}/\texttt{enwik9} is predominantly retrieval, not
high-order statistics~\cite{paq}; Steinr\"ucken measured the Bayesian
version to be competitive with its ad hoc
counterparts~\cite{steinruecken}.  In-house, by elimination: the
pooled-lag run found the entire tail of longer lags worth $0.044$
bits per token --- long-range \emph{pairwise} statistics are nearly
empty --- yet the files demonstrably repeat over exactly those
ranges.  What remains at long range is retrievable, not statistical,
and nothing in our current machinery can express it.  This is the one
route whose ceiling we have not measured even indirectly.

One more piece of evidence matters, because the long-term goal of
this program is a model of $p(\text{current} \mid \text{past})$ in
the sense of the language-model literature, not a file compressor:
copying is demonstrably a core mechanism of that world too.  When
transformer language models are opened up, a large part of their
in-context ability is carried by ``induction heads'' --- circuits
that find the previous occurrence of the current context and copy
what followed it, which is this route's latent pointer implemented in
attention~\cite{inductionheads}.  And adding an explicit retrieval
channel to a trained language model --- predicting by interpolating
the model with a distribution over what followed the nearest matching
contexts in a datastore --- lowers perplexity
directly~\cite{knnlm}.  So the pointer is not a compression trick
that would be discarded on the way to the real goal; it is a
mechanism the real goal measurably uses.

\paragraph{Route 4 --- Switching: a mixture over sequences of members.}

\emph{Prediction.}  Once more we stand at position $t$.  Each member
$k$ of a family already supplies a sequential prediction
$p_k(\,\cdot \mid x^{t-1})$; nothing about the members changes.  The
family average used throughout picks the member once, for the whole
file: it costs $\log_2 K$ bits for $K$ members, and its posterior is
consistently degenerate (\texttt{posterior\_max} $= 1.0$ in the
pairwise runs) --- one member wins globally.  But a global winner can
hide local losers.  The parked observation that one first-order state
cannot serve a stream of two interleaved languages is exactly this.

The repair is to let the active member change along the file.  The
latent variable becomes a piecewise-constant sequence of members: at
each position the active member persists, or, with a small switch
probability $\alpha_t$ (say $\alpha_t \propto 1/t$, so that
logarithmically many switches are expected a priori), is redrawn from
the prior over members.  The decoder keeps one weight per member ---
the posterior probability that member $k$ is active now --- updates
the weights by the forward recursion (multiply each by that member's
probability for the symbol just decoded, then move a fraction
$\alpha_t$ of the total mass to the prior), and predicts by the
weighted average.  This is the switch distribution~\cite{switching}:
exact Bayes over segmentations, $O(K)$ per position, no weight
learned.

\emph{The codelength.}  As in Route 3, the posterior is
path-dependent and there is no closed form; the codelength is
accumulated by the sequential walk, the switch layer adding $O(K)$
arithmetic per position on top of the members' own predictions.  At
$\alpha \equiv 0$ the scheme reduces exactly to the current family
code.  In general it pays about $\log_2 n$ bits per switch actually
used: it tracks the best \emph{segmentation} of the file into member
regimes, not the best member.

\emph{Fit with our scheme.}  Drop-in: the members are the members we
already have --- and, later, routes 1, 2 and 3 as members beside
them.
The interface is each member's sequential prediction, which every
member already exposes.

\emph{Cost.}  $O(K)$ per position and $O(K)$ memory on top of the
members' own predictions.  The expensive part is the members, not the
switch.

\emph{Evidence.}  Published: the switch distribution resolves the
catch-up phenomenon that makes single-model Bayes slow to change
regimes~\cite{switching}, with regret logarithmic per switch; Volf
and Willems measured that switching between CTW and a match-model
companion beat both alone~\cite{volf} --- the historical composition
of routes 1 and 3 under route 4.  In-house: the degenerate family
posteriors, and the two-languages observation, are both direct
symptoms of the missing latent axis.

\paragraph{Order of attack.}  The measured numbers, and the state of
the code, rank the routes.  Constraint models first: they are the
running experiment continued, their two-lag core is already
implemented and measured, and each extension --- more lags, lag--lag
constraints, triples --- reuses machinery that exists today.
Hierarchy next: it attacks the $0.400$ bits per character of learning
shortfall head-on and composes with the evaluator with no sequential
walk, but its reach in distance is short --- a tree over a handful of
recent symbols is still a Markov model of modest span, however
adaptively pruned.  Retrieval third in order but largest in ceiling:
it is the only route that reaches arbitrary distances, the only one
cheap at the full vocabulary immediately, and the one the
language-model evidence singles out.  Switching last: the glue,
principled and nearly free, but it amplifies the others rather than
saving bits by itself.  The first two extend memory in \emph{order},
the third in \emph{distance}; they are complements, not rivals.  All
four are the move this paper is built on --- a latent variable, a
prior, an exact algorithm --- so adopting them changes the program's
scale, not its character.

\begin{thebibliography}{9}
\bibitem{layered2026}
\emph{A Layered Simplex Architecture for Large Alphabets}, 2026.
Companion paper and repository (\texttt{product\_model}).
\bibitem{ctw95}
F.~M.~J. Willems, Y.~M. Shtarkov, T.~J. Tjalkens.
The context-tree weighting method: basic properties.
\emph{IEEE Trans.\ Inform.\ Theory} 41(3):653--664, 1995.
\bibitem{memoizer}
F.~Wood, J.~Gasthaus, C.~Archambeau, L.~James, Y.~W. Teh.
The sequence memoizer.
\emph{Communications of the ACM} 54(2):91--98, 2011.
\bibitem{ppmstar}
J.~G. Cleary, W.~J. Teahan.
Unbounded length contexts for PPM.
\emph{The Computer Journal} 40(2/3):67--75, 1997.
\bibitem{steinruecken}
C.~Steinr\"ucken.
\emph{Lossless Data Compression}.
PhD thesis, University of Cambridge, 2014.
\bibitem{paq}
M.~V. Mahoney.
Adaptive weighing of context models for lossless data compression.
Florida Tech.\ Technical Report CS-2005-16, 2005.
\bibitem{switching}
T.~van Erven, P.~Gr\"unwald, S.~de Rooij.
Catching up faster by switching sooner: a predictive approach to
adaptive estimation.
\emph{J.\ Royal Statist.\ Soc.\ B} 74(3):361--417, 2012.
\bibitem{volf}
P.~A.~J. Volf, F.~M.~J. Willems.
Switching between two universal source coding algorithms.
\emph{Proc.\ Data Compression Conference}, 1998.
\bibitem{rosenfeld}
R.~Rosenfeld.
A maximum entropy approach to adaptive statistical language
modeling.
\emph{Computer Speech and Language} 10(3):187--228, 1996.
\bibitem{yedidia}
J.~S. Yedidia, W.~T. Freeman, Y.~Weiss.
Constructing free-energy approximations and generalized belief
propagation algorithms.
\emph{IEEE Trans.\ Inform.\ Theory} 51(7):2282--2312, 2005.
\bibitem{inductionheads}
C.~Olsson et al.
In-context learning and induction heads.
\emph{Transformer Circuits Thread}, Anthropic, 2022.
\bibitem{knnlm}
U.~Khandelwal, O.~Levy, D.~Jurafsky, L.~Zettlemoyer, M.~Lewis.
Generalization through memorization: nearest neighbor language
models.
\emph{Proc.\ ICLR}, 2020.
\end{thebibliography}

\end{document}
